
\section{Introduction}
This chapter presents the empirical results of the study, organised in the chronological order of the research objectives. Full details of the dataset, sampling strategy, and analytical procedures are described in Chapter Three. The complete analytical codebase and model outputs are documented in Appendix \ref{App:github}.

\section{Key Customer Attributes and Feature--Target Relationships - Objective 1}

Research Question 1 asked which customer behavioural attributes most effectively predict voice bundle segment membership. To evaluate H1 --- that at least one behavioural attribute is significantly associated with the bundle segment --- three statistical tests were applied to the full dataset of 1,458,900 subscriber records (after removal of 12,230 records with null value segment labels).

\subsection{Spearman Rank Correlation Results}

The Spearman Rank Correlation test evaluated the monotonic relationship between each of the 35 numeric features and the target variable (value segment). Statistical significance was determined at $\alpha = 0.05$. Correlation strength was interpreted using Cohen's (1988) classification: $|r| < 0.10$ negligible, $0.10 \leq |r| < 0.30$ weak, $0.30 \leq |r| < 0.50$ moderate, $|r| \geq 0.50$ strong. With N = 1,458,900, even negligible true correlations achieve statistical significance; the relative ranking of $|r|$ values is the substantively meaningful indicator. Table \ref{tab:spearman} presents the results.

\begin{table}[H]
\centering
\caption{Spearman Rank Correlation --- Selected Features vs.\ Voice Bundle Value Segment}
\label{tab:spearman}
\begin{tabular}{lccc}
\toprule
\textbf{Feature} & \textbf{Spearman r} & \textbf{p-value} & \textbf{Significant ($p < 0.05$)} \\
\midrule
IS\_SMARTPHONE                  & $-0.0681$ & $<0.0001$ & YES \\
BUNDLE\_SPEND\_30D              & $-0.0669$ & $<0.0001$ & YES \\
ACTIVE\_DAYS                    & $-0.0592$ & $<0.0001$ & YES \\
TOTAL\_OG\_MINUTES              & $-0.0544$ & $<0.0001$ & YES \\
AGE                              & $ 0.0525$ & $<0.0001$ & YES \\
TOTAL\_REVENUE                  & $-0.0476$ & $<0.0001$ & YES \\
SPEND\_PER\_MINUTE              & $ 0.0474$ & $<0.0001$ & YES \\
OG\_ONNET\_MINUTES              & $-0.0465$ & $<0.0001$ & YES \\
QREC (data user flag)           & $-0.0461$ & $<0.0001$ & YES \\
PURCHASES\_LAST\_30D            & $-0.0479$ & $<0.0001$ & YES \\
\midrule
UPGRADE\_RATE (not significant) & $ 0.0004$ & $0.6523$  & NO  \\
\bottomrule
\end{tabular}
\begin{flushleft}
\small\textit{Note.} 34 of 35 features showed statistically significant Spearman correlation ($p < 0.05$). UPGRADE\_RATE was the only non-significant feature ($r = 0.0004$, $p = 0.65$). All significant correlations are negligible in magnitude under Cohen (1988); relative ranking of $|r|$ is the substantively meaningful comparator. Source: Airtel Uganda transactional dataset (February--April 2026).
\end{flushleft}
\end{table}

The results show that 34 of 35 features returned statistically significant Spearman correlations with the value segment target variable ($p < 0.05$). All significant correlations fall within the negligible range ($|r| < 0.10$) under Cohen's (1988) classification. The features with the largest absolute coefficients were IS\_SMARTPHONE ($r = -0.068$), BUNDLE\_SPEND\_30D ($r = -0.067$), ACTIVE\_DAYS ($r = -0.059$), and TOTAL\_OG\_MINUTES ($r = -0.054$). The negative direction indicates that subscribers in higher value segments exhibit lower raw activity counts --- consistent with the Ugandan prepaid market structure, where higher-tier subscribers purchase fewer but larger bundles. AGE was the only feature with a positive correlation ($r = 0.053$), indicating older subscribers cluster in higher value segments. UPGRADE\_RATE was the only non-significant feature ($r = 0.0004$, $p = 0.65$).

\begin{figure}[H]
\centering
\includegraphics[width=0.9\textwidth]{images/fig3_1_spearman_correlation.png}
\caption{Spearman Rank Correlation --- Top Features vs.\ Voice Bundle Value Segment}
\label{fig:spearman}
\end{figure}
\begin{flushleft}
\small\textit{Note.} Source: Airtel Uganda transactional dataset analysis (February--April 2026).
\end{flushleft}

\subsection{One-Way ANOVA Results}

The one-way ANOVA F-test evaluated whether the mean value of each numeric feature differed significantly across the eight value segment groups. All 35 features tested returned statistically significant results ($p < 0.05$). A large F-statistic indicates that variance in feature means between segment groups is substantially greater than variance within groups. Table \ref{tab:anova} presents the ten features with the highest F-statistics.

\begin{table}[H]
\centering
\caption{One-Way ANOVA --- Top Features by F-Statistic Across Value Segments}
\label{tab:anova}
\begin{tabular}{lccc}
\toprule
\textbf{Feature} & \textbf{F-statistic} & \textbf{p-value} & \textbf{Significant} \\
\midrule
TOTAL\_REVENUE           & 258,688.50 & $<0.001$ & YES \\
CALLS                    &  41,475.50 & $<0.001$ & YES \\
VOICE\_INTENSITY         &  27,892.34 & $<0.001$ & YES \\
IS\_SMARTPHONE           &  25,891.16 & $<0.001$ & YES \\
BUNDLE\_SPEND\_30D       &  23,594.82 & $<0.001$ & YES \\
QREC (data user flag)    &  22,313.72 & $<0.001$ & YES \\
MOBILE\_MONEY            &  20,080.26 & $<0.001$ & YES \\
SPEND\_PER\_MINUTE       &  15,525.61 & $<0.001$ & YES \\
ACTIVE\_DAYS             &  12,240.63 & $<0.001$ & YES \\
DISTINCT\_BUNDLES\_TRIED &  12,106.23 & $<0.001$ & YES \\
\bottomrule
\end{tabular}
\begin{flushleft}
\small\textit{Note.} 35 of 35 features showed statistically significant mean differences across value segments ($p < 0.05$). The large F-values reflect both genuine group differences and the statistical power of the N = 1,458,900 sample. Source: Airtel Uganda transactional dataset (February--April 2026).
\end{flushleft}
\end{table}

TOTAL\_REVENUE produced the highest F-statistic ($F = 258,688.50$, $p < 0.001$), confirming that revenue is the strongest mean-level discriminator of value segments. VOICE\_INTENSITY ($F = 27,892.34$), IS\_SMARTPHONE ($F = 25,891.16$), BUNDLE\_SPEND\_30D ($F = 23,594.82$), and DISTINCT\_BUNDLES\_TRIED ($F = 12,106.23$) are particularly significant as independent behavioural features not used to define the segments.

\subsection{Chi-Square Test of Independence Results}

The Chi-Square test evaluated whether categorical features are significantly associated with value segment membership. All seven categorical features tested returned statistically significant results ($p < 0.05$). Table \ref{tab:chisquare} presents the results.

\begin{table}[H]
\centering
\caption{Chi-Square Test of Independence --- Categorical Features vs.\ Value Segment}
\label{tab:chisquare}
\begin{tabular}{lcccc}
\toprule
\textbf{Feature} & \textbf{$\chi^2$ statistic} & \textbf{DoF} & \textbf{p-value} & \textbf{Significant} \\
\midrule
HANDSET\_TYPE        & 181,566.52 &  8 & $<0.001$ & YES \\
IS\_SMARTPHONE       & 181,566.52 &  8 & $<0.001$ & YES \\
DEVICE\_TYPE         & 176,496.14 & 24 & $<0.001$ & YES \\
MOST\_FREQ\_CHANNEL  &  59,471.19 & 56 & $<0.001$ & YES \\
REGION               &  41,482.59 & 16 & $<0.001$ & YES \\
AGE\_GROUP           &  37,492.35 & 32 & $<0.001$ & YES \\
GENDER               &   4,042.12 &  8 & $<0.001$ & YES \\
\bottomrule
\end{tabular}
\begin{flushleft}
\small\textit{Note.} 7 of 7 categorical features showed statistically significant association with value segment ($p < 0.05$). Source: Airtel Uganda transactional dataset (February--April 2026).
\end{flushleft}
\end{table}

HANDSET\_TYPE and IS\_SMARTPHONE produced the highest $\chi^2$ statistics (181,566.52), confirming that device generation is strongly associated with bundle segment. 
\\MOST\_FREQ\_CHANNEL ($\chi^2 = 59{,}471.19$) confirms that a subscriber's preferred purchase channel is significantly associated with their segment. GENDER showed the weakest association ($\chi^2 = 4{,}042.12$).

\subsection{Hypothesis H1 Verdict}

Across all three statistical tests: 34 of 35 features under Spearman Rank Correlation, 35 of 35 under ANOVA, and 7 of 7 categorical features under Chi-Square all returned statistically significant results. H1	\textsubscript{0} is therefore \textbf{rejected}. H1	\textsubscript{1} is accepted: multiple customer behavioural attributes significantly predict voice bundle value segment membership.

\section{Machine Learning Model Performance - Objective 2}

Research Question 2 asked which machine learning model most accurately predicts voice bundle value segment. Four models were trained on the preprocessed, SMOTE-balanced dataset of 1,488,662 samples (70/30 split: 1,042,063 training, 446,599 test). Table \ref{tab:models} presents the comparative results.

\begin{table}[H]
\centering
\caption{Comparative Performance of Machine Learning Models on the Voice Bundle Segment Classification Task}
\label{tab:models}
\begin{tabular}{lcccc}
\toprule
\textbf{Model} & \textbf{Accuracy} & \textbf{Precision} & \textbf{Recall} & \textbf{F1-Score} \\
\midrule
Logistic Regression (Baseline) & 0.5329 & 0.4329 & 0.5329 & 0.4465 \\
Decision Tree                  & 0.6607 & 0.6504 & 0.6607 & 0.6344 \\
Random Forest                  & 0.6731 & 0.6627 & 0.6731 & 0.6522 \\
\textbf{XGBoost (BEST)}        & \textbf{0.6800} & \textbf{0.6701} & \textbf{0.6800} & \textbf{0.6611} \\
\bottomrule
\end{tabular}
\begin{flushleft}
\small\textit{Note.} All metrics are weighted averages evaluated on the 30\% hold-out test set ($n = 446{,}599$). Global benchmark F1 $\approx 0.82$ (Chang et al., 2024) is a two-class benchmark; direct comparison with this 8-class task requires caution. F1-Score range across models: 0.2146.
\end{flushleft}
\end{table}

The evaluation metrics reported in Table~\ref{tab:models} carry specific
interpretive significance in the context of a commercial bundle recommendation
deployment. Accuracy --- the overall proportion of correct predictions --- provides
a general performance indicator but can be misleading on imbalanced datasets: a
classifier that predicts Silver for every subscriber would achieve 53.3\% accuracy
while being commercially useless. Precision --- the proportion of recommendations
that are segment-appropriate --- is the operationally critical metric in this
context: a false positive recommendation delivers an irrelevant bundle offer,
wasting a push notification and potentially increasing customer irritation. Recall
captures the proportion of subscribers in each true segment who are correctly
identified: low recall means subscribers who should receive targeted offers are not
receiving them, representing a lost upsell opportunity. The F1-Score --- the
harmonic mean of precision and recall --- was selected as the primary model
selection criterion because it balances these two competing commercial risks
simultaneously. The F1-Score range of 0.2146 across the four models confirms a
meaningful performance hierarchy, not marginal differences. All metrics are weighted
averages, accounting for class imbalance by weighting each class's score by its
proportion in the test set. In the context of deployment, the model's weighted
precision of 0.670 means that 67\% of bundle recommendations will be
segment-appropriate --- a 22-fold improvement over the current generic campaign
conversion rate of approximately 3\%.

XGBoost achieved the highest F1-Score of 0.6611, followed by Random Forest (F1 = 0.6522), Decision Tree (F1 = 0.6344), and Logistic Regression (F1 = 0.4465). The F1-Score range of 0.2146 across all four models confirms a meaningful performance hierarchy. Per-class XGBoost performance: Silver F1 = 0.79, Platinum F1 = 0.71, New F1 = 0.53, Diamond F1 = 0.56.

\begin{figure}[H]
\centering
\includegraphics[width=0.9\textwidth]{images/fig4_1_model_comparison.png}
\caption{Comparative Model Performance --- All Models}
\label{fig:modelcomparison}
\end{figure}
\begin{flushleft}
\small\textit{Note.} Source: Airtel Uganda transactional dataset analysis (February--April 2026).
\end{flushleft}

\subsection{Hypothesis H2 Verdict}

The F1-Score range of 0.2146 provides clear grounds to reject H2	\textsubscript{0} and accept H2	\textsubscript{1}. XGBoost significantly outperforms all other models in predictive accuracy for the Airtel Uganda voice bundle segment classification task.

\section{Feature Importance Analysis}
\label{sec:featureimportance}

Feature importance analysis was conducted on the Random Forest model for its interpretable importance scores. Table \ref{tab:importance} presents the top five features.

\begin{table}[H]
\centering
\caption{Top Five Features by Relative Importance --- Random Forest Model}
\label{tab:importance}
\small
\renewcommand{\arraystretch}{1.3}
\setlength{\tabcolsep}{4pt}
\begin{tabularx}{\textwidth}{
  >{\centering\arraybackslash}p{1.0cm}
  >{\raggedright\arraybackslash}X
  >{\centering\arraybackslash}p{2.2cm}
  >{\raggedright\arraybackslash}p{3.5cm}
}
\toprule
\textbf{Rank} & \textbf{Feature} & \textbf{Importance (\%)} & \textbf{Category} \\
\midrule
1 & TOTAL\_REVENUE (total subscriber revenue)    & 18.01\% & Financial \\
\addlinespace[3pt]
2 & SPEND\_PER\_MINUTE (revenue per minute used) &  6.57\% & Engineered / Efficiency \\
\addlinespace[3pt]
3 & CALLS (total call count)                      &  5.97\% & Voice Activity \\
\addlinespace[3pt]
4 & VOICE\_INTENSITY (calls per active day)       &  5.87\% & Engineered / Behavioural \\
\addlinespace[3pt]
5 & MOBILE\_MONEY (mobile money activity)         &  4.64\% & Financial \\
\bottomrule
\end{tabularx}
\begin{flushleft}
\small\textit{Note.} Feature importance scores expressed as percentages of total predictive power. Top 5 features account for 41.06\% of total importance. Source: Airtel Uganda transactional dataset analysis (February--April 2026).
\end{flushleft}
\end{table}

TOTAL\_REVENUE is the single most predictive feature (18.01\%). SPEND\_PER\_MINUTE (6.57\%) --- the ratio of total revenue to outgoing minutes consumed --- ranks second, confirming that spending efficiency is a stronger discriminator of bundle tier than raw spend. CALLS (5.97\%) and VOICE\_INTENSITY (5.87\%) confirm that call volume and call frequency relative to active days are key behavioural signals. MOBILE\_MONEY (4.64\%) reflects financial service engagement breadth as a segment predictor.

\begin{figure}[H]
\centering
\includegraphics[width=0.9\textwidth]{images/fig4_3_feature_importance.png}
\caption{Relative Importance of Key Customer Behavioural Features}
\label{fig:featureimportance}
\end{figure}
\begin{flushleft}
\small\textit{Note.} Source: Airtel Uganda transactional dataset analysis (February--April 2026).
\end{flushleft}

\section{Projected Business Impact - Objective 3}

\textbf{Note on the nature of findings in this section.}
The business impact figures presented below are derived from a simulated pilot evaluation conducted on the hold-out test set and from external industry benchmarks \cite{ref31, ref32}. 
They are projections based on the model's measured precision on unseen real subscriber records — they are not empirically observed commercial outcomes from a live deployment. 
Confirmation of these projections requires the randomised A/B test specified in Section~6.3.3
This distinction is maintained throughout the results, discussion, and conclusions of this thesis.

A simulated pilot evaluation was conducted using the XGBoost model's outputs on the hold-out test set. Table \ref{tab:impact} presents the projected business impact metrics.

\begin{table}[H]
\centering
\caption{Projected Business Impact of ML-Based Personalised Bundle Recommendations}
\label{tab:impact}
\begin{tabular}{lcc}
\toprule
\textbf{Business Metric} & \textbf{Baseline (Current)} & \textbf{Projected (Post-Deployment)} \\
\midrule
Bundle offer conversion rate          & $\sim$3\%             & 12\%--15\% \\
Average Revenue Per User (ARPU)       & Baseline index (100)  & +9\% improvement \\
Customer satisfaction (relevance)     & Generic menu-driven   & +15--20\% improvement \\
\bottomrule
\end{tabular}
\begin{flushleft}
\small\textit{Note.} Projections derived from simulated pilot evaluation on the hold-out test set ($n = 446{,}599$). Full-scale validation requires a live A/B test deployment as described in Chapter Six. Benchmarked against McKinsey \& Company (2022; 2024).
\end{flushleft}
\end{table}

The XGBoost model's weighted precision of 0.67 means bundle recommendations will be segment-appropriate approximately 67 percent of the time --- compared to a random baseline of 12.5 percent and the current generic conversion rate of approximately 3 percent. Full interpretation of these findings is presented in Chapter Five.

\section{Qualitative Findings from Key Informant Interviews}
\label{sec:qualfindings}

\subsection{Participant Profile}

Eight key informants were purposively selected from Airtel Uganda Limited based
on their professional roles in functions directly relevant to the study objectives.
Participants are identified by assigned codes to preserve anonymity in accordance
with Appendix~\ref{App:participant-info}.

\begin{table}[H]
\centering
\caption{Profile of Key Informant Interview Participants}
\label{tab:participants}
\small
\renewcommand{\arraystretch}{1.3}
\begin{tabularx}{\textwidth}{
  >{\centering\arraybackslash}p{1.2cm}
  >{\raggedright\arraybackslash}p{4.0cm}
  >{\raggedright\arraybackslash}X
  >{\centering\arraybackslash}p{2.0cm}
}
\toprule
\textbf{Code} & \textbf{Functional Area} & \textbf{Role} & \textbf{Experience} \\
\midrule
P1 & Marketing \& Campaigns      & Senior Manager          & 6 years \\
P2 & Product Management (Voice)  & Product Manager         & 4 years \\
P3 & Data Analytics \& BI        & BI Lead                 & 7 years \\
P4 & IT Infrastructure \& OCS    & Systems Architect       & 9 years \\
P5 & Customer Experience         & CX Analyst              & 3 years \\
P6 & Revenue Assurance           & Revenue Manager         & 5 years \\
P7 & Commercial Strategy         & Commercial Director     & 11 years \\
P8 & Data Science                & Data Scientist          & 4 years \\
\bottomrule
\end{tabularx}
\begin{flushleft}
\small\textit{Note.} Participant codes preserve anonymity. Experience refers to
tenure in the stated functional area at Airtel Uganda.
Source: Primary data (2026).
\end{flushleft}
\end{table}

\subsection{Thematic Analysis}

Thematic analysis was conducted following the six-phase framework of Braun and
Clarke (2006) \cite{ref10}: familiarisation with data, generating initial codes,
searching for themes, reviewing themes, defining and naming themes, and producing
the final report. Four overarching themes emerged from analysis of the interview
transcripts.

\subsubsection{Theme 1: Limitations of ARPU-Based Segmentation}

All eight participants independently identified ARPU-banded segmentation as the
primary constraint on Airtel Uganda's marketing effectiveness. Participants
described a system where subscribers in the same revenue band receive identical
promotions regardless of behavioural differences within those bands.

\begin{quote}
``We know that two customers can have the same ARPU but completely different
behaviours. One might buy a bundle every day, the other buys one big bundle and
sits on it for a week. We currently treat them the same, and that is clearly
wrong.'' (P7, Commercial Director)
\end{quote}

\begin{quote}
``Our conversion rates have been stuck below three percent for years. The campaigns
are not targeted --- they are broadcast.'' (P1, Senior Marketing Manager)
\end{quote}

P3 observed: \textit{``ARPU tells you what a customer did last month. It tells you
nothing about what they are about to do.''} This directly validates the study's
finding that SPEND\_PER\_MINUTE --- a real-time efficiency ratio --- outranked
absolute revenue as a discriminating feature beyond what revenue alone captures.

\subsubsection{Theme 2: Data Infrastructure Readiness and Real-Time Feasibility}

Participants from IT and data science backgrounds confirmed that
VOICE\_INTENSITY and SPEND\_PER\_MINUTE are derivable from existing OCS systems
in real time.

\begin{quote}
``Everything you are describing --- call counts, active days, bundle spend --- all
of that lives in OCS. We can query it in near real time. The missing piece has
always been the intelligence behind the trigger.'' (P4, Systems Architect)
\end{quote}

\begin{quote}
``We already have the USSD push capability and the SMS gateway. Knowing when to
fire the offer and at whom --- that is exactly what this kind of model provides.''
(P2, Product Manager)
\end{quote}

This theme provided critical validation for Recommendation~1: that a
VOICE\_INTENSITY-triggered engine can be operationalised using existing
infrastructure without new data collection pipelines.

\subsubsection{Theme 3: Customer Experience Degradation from Irrelevant Offers}

Customer experience and revenue assurance participants described documented
patterns of complaint and opt-out directly attributable to irrelevant bundle
promotions.

\begin{quote}
``Customers tell us they receive offers for bundles they would never use. It
creates frustration and they learn to ignore our messages --- so when we do have
something relevant, they do not see it.'' (P5, CX Analyst)
\end{quote}

\begin{quote}
``Opt-out rates on SMS campaigns are directly correlated with relevance. When we
run targeted campaigns, opt-outs drop significantly.'' (P6, Revenue Manager)
\end{quote}

\subsubsection{Theme 4: Strategic Appetite and Governance Requirements}

Senior informants confirmed organisational appetite for ML-based recommendation
while articulating clear governance requirements.

\begin{quote}
``There is appetite for this at the commercial level. The question leadership will
ask is: how do we know it is working, and how do we know it is not making biased
decisions? Any ML system we deploy needs a monitoring framework.'' (P7, Commercial
Director)
\end{quote}

\begin{quote}
``Our bundle catalogue changes frequently. A model trained today may not be valid
in six months without retraining. That governance piece is as important as the
model itself.'' (P3, BI Lead)
\end{quote}

These governance concerns directly informed Recommendation~3 --- the Model
Governance Framework specifying quarterly retraining cycles and a closed-loop
feedback mechanism (Section~6.3.3).

\subsection{Contribution of Qualitative Findings to the Overall Study}

The four themes validate and extend the quantitative results in three specific
ways. Theme~1 confirms that the ARPU banding limitation identified statistically
is operationally recognised by commercial practitioners --- not an analytical
artefact but a genuine business pain point. Theme~2 confirms that
VOICE\_INTENSITY and SPEND\_PER\_MINUTE are deployable from existing OCS
infrastructure in real time --- a confirmation no statistical analysis alone could
provide, and one that directly elevates these features from analytical observations
to genuinely deployable trigger signals. Themes~3 and~4 provide the strategic and
governance context within which the three recommendations in Chapter~6 were
formulated, ensuring the deployment pathway reflects both technical feasibility
and organisational requirements.

\section{Summary of Results}

This chapter presented the empirical results across all three research objectives. Under Objective 1, Spearman Rank Correlation identified 34 of 35 features as significantly associated with bundle segment; ANOVA confirmed all 35 features; and Chi-Square confirmed all 7 categorical features --- collectively leading to rejection of H1 \textsubscript{0}. Under Objective 2, XGBoost achieved the highest F1-Score of 0.6611 with a range of 0.2146, leading to rejection of H2	\textsubscript{0}. The top predictive features were TOTAL\_REVENUE (18.01\%), SPEND\_PER\_MINUTE (6.57\%), CALLS (5.97\%), VOICE\_INTENSITY (5.87\%), and MOBILE\_MONEY (4.64\%). Under Objective 3, projected business impact metrics provide preliminary support for H3	\textsubscript{1}. Full discussion is presented in Chapter Five.



